\section{Introduction}\label{sec:intro}

Sensing devices that people wear or carry, such as watches, phones and fitness bands, generate continuous streams of inertial data that are useful for health monitoring, rehabilitation and context-aware services, and that are too personal and too voluminous to be collected centrally. Federated learning (FL) trains models across such devices without moving the raw data \citep{mcmahan2017communication,kairouz2021advances}. In practice, however, an FL round is not free for the device: local training draws on a battery that must also serve the user, and the model that comes out of training will run on the same device, around the clock, for as long as it is deployed. Energy is therefore the scarce resource of cross-device learning at the edge, and it enters twice---once when the model is trained and again, every second, when it is used.

The standard FL toolbox treats energy indirectly. Client-selection rules choose participants by statistical utility and speed \citep{lai2021oort}; compression reduces the bits sent \citep{konecny2016federated,alistarh2017qsgd}; capacity-aware methods give weak devices a smaller sub-network \citep{diao2021heterofl,kim2023depthfl,ilhan2023scalefl}; wireless-FL studies optimise transmit power and CPU frequency for a fixed model \citep{tran2019federated,yang2021energy}. Each acts on one knob, usually with a fixed assignment, and none of them asks what a Joule of training buys, how that answer differs between a fast, efficient device and a slow, inefficient one, or whether the cheaper updates still give a model that is trained everywhere it needs to be. On the deployment side, early-exit networks reduce inference energy by letting easy inputs leave early \citep{teerapittayanon2016branchynet,kaya2019shallow}, but their training in federated, battery-limited fleets is rarely considered together with the energy it consumes.

\paragraph{The idea}
A multi-exit network offers a natural handle on energy: a device that trains only the first $d$ blocks pays roughly in proportion to $d$, and the shallow blocks are the ones that every exit relies on. The risk is that the deep blocks, and the deep exits, are then trained too rarely. Our starting point is to treat this risk as a precisely quantifiable bias (Theorem~\ref{thm:main}): when selected clients train block $j$ with probability $\sigma_j$, the aggregate update is biased by an amount that grows with $1-\sigma_j$ and produces an error floor that no number of rounds removes. This turns a heuristic worry into a design constraint---keep $\sigma_j$ above a target---and the constraint can be enforced online by a queue.

\paragraph{EcoFed}
We build a controller, EcoFed, around that constraint. Each round, every client scores its feasible configurations (depth, upload precision, local epochs) by a utility, an energy cost priced by a per-client queue and a shadow price, and a bonus for reaching blocks whose coverage queue is backlogged; the server then picks the best few. Energy queues measure spending against a per-client budget in units of the budget itself, so weak, small-battery devices are priced out of expensive configurations without being excluded. Coverage queues keep deep blocks trained. The controller needs no model of the learning dynamics beyond a monotone utility and uses the energy that the device reports, so it keeps working when its own energy model is wrong (Section~\ref{sec:sens}).

\paragraph{What we find}
We evaluate on two public wearable benchmarks with natural person-level non-IID partitions, comparing against seven baselines (FedAvg, FedProx, quantised FedAvg, a paced FedAvg, static-depth heterogeneity, an Oort-style selector and an energy-greedy selector) under three battery tightnesses, with tuning on separate seeds and validation subjects only. The quantitative findings are reported in Section~\ref{sec:results} and summarised in the abstract. We also analyse the lifecycle of the trained model, including the crossover point at which serving the model costs more energy than training it did, which shows how much room inference-side savings leave for training-side choices.

\paragraph{Scope and honesty about energy}
All energy figures in this paper are \emph{model-based}: constants for three device classes are taken from typical published orders of magnitude, and we did not have a power meter. We therefore treat robustness to energy-model error as a first-class experiment, and we state the consequence plainly: the paper establishes how a controller should behave given an energy model and that it degrades gracefully when the model is wrong, not the absolute number of Joules on any particular board.

\paragraph{Contributions}
\begin{enumerate}[leftmargin=1.6em,itemsep=1pt]
\item A device-aware energy model for federated multi-exit training and inference that separates arithmetic, memory traffic and radio energy, together with a lifecycle accounting that includes a break-even and a training--inference crossover (Section~\ref{sec:model}).
\item A convergence bound for federated training with heterogeneous client depth in which skipped exits create a coverage-dependent error floor, and a corollary that gives the coverage needed for a target accuracy (Section~\ref{sec:theory}, \ref{app:proof}).
\item EcoFed, an online drift-plus-penalty controller that jointly chooses clients, depth, precision and epochs under per-client energy queues, a shadow energy price and coverage queues, with $O(N|\mathcal C|)$ cost per round (Section~\ref{sec:controller}).
\item An evaluation on two wearable benchmarks with paired statistics, ablations, a test of the coverage prediction, robustness to energy-model error and to the radio-to-compute cost ratio, fairness and lifetime, and a lifecycle analysis (Sections~\ref{sec:setup}--\ref{sec:results}).
\end{enumerate}
The remainder of the paper is organised as follows. Section~\ref{sec:related} reviews related work, Section~\ref{sec:model} defines the model, Section~\ref{sec:method} presents theory and algorithm, Section~\ref{sec:setup} the experimental setup, Section~\ref{sec:results} the results, Section~\ref{sec:discussion} discusses implications and limitations.
